\documentclass[a4paper,12pt]{report}
\usepackage[italian]{babel}
\usepackage[latin1]{inputenc}
\usepackage{graphicx}
\setcounter{tocdepth}{2}
\begin{document}

%\begin{figure}
%\centering
%\includegraphics[scale=0.15]{Logo.jpg}
%\end{figure}
\begin{center}
\textbf{UNIVERSIT\'{A} DEGLI STUDI "ROMA TRE"}\\
\vspace{0.3cm}
\textmd{\large Dipartimento di Matematica e Fisica}\\
\vspace{0.2cm}
\large{ Corso di Laurea Magistrale in Scienze Computazionali}\\	
\vspace{1.4cm}	
\textmd{\Large{Tesi di Laurea Magistrale}}\\
\vspace{1.6cm}
\LARGE{\LARGE{\textbf{Ricerca della tolopologia ottimale di un sistema di deep learning per identificazioni di oggetti architettonici}}}\\
\vspace{5cm}
\begin{tabular}{ccccccccccc}
	\large{Candidato}& & & & & & & & & & \large{Relatore}\\
	\large{D�sir�e Adiutori} & & & & & & & & & &\large{Prof. Alberto Paoluzzi}
\end{tabular}\\
\vspace{4cm}
\normalsize{Anno Accademico 2017/2018}\\
\normalsize{Luglio 2018}
\end{center}
\newpage
\tableofcontents
\newpage
\section*{Introduzione}
\addcontentsline{toc}{section}{Introduzione}
Dalla scoperta dei computer, l'uomo fa sempre pi� affidamento alle macchine per risolvere problemi complessi di calcolo. Con l'aumentare delle prestazioni dei computer, man mano si sono sviluppati algoritmi di calcolo sempre pi� efficienti.
Nel 1959 l'ingegnere del MIT, Arthur Samuel coni� il termine \textit{"machine learning"}, descrivendo l'apprendimento automatico come un "campo di studio che d� ai computer la possibilit� di apprendere senza essere programmati esplicitamente per farlo".\\
Definiamo l'apprendimento automatico come un insieme di metodi in grado di rilevare automaticamente i modelli nei dati e quindi utilizzare i modelli scoperti per prevedere i dati futuri o per eseguire altri tipi di processi decisionali in condizioni di incertezza. L'insegnamento alla macchine �, pertanto, imprescindibile dai dati. Pi� dati si passano alla macchina, pi� pu� imparare. Per questo motivo con l'avvento di Internet, dagli anni '90 ad oggi il tema del "machine learning" � diventato sempre pi� attuale, la mole di dati reperibile dal web � cospicuo e ha permesso che questo campo sia espnenzialmente progredito.\\
Il \textit{"deep learning"} � un tipo particolare di machine learning che riguarda l'emulazione di come gli esseri umani apprendono. Esso affronta i problemi del machine learning, rappresentando il mondo come una gerarchia di concetti annidati, dove ogni concetto � definito in relazione a concetti pi� semplici, e le rappresentazioni astratte vengono calcolate in termini di concetti meno astratti. Il Deep Learning implica l'utilizzo di reti neurali artificiali (\textit{deep artificial neural networks}), algoritmi e sistemi computazionali ispirati dal cervello umano, per affrontare i problemi del Machine Learning. L'analogia di Shehzad Noor Taus Priyo pu� aiutare a capire meglio cosa siano le reti neurali:\\
"Immaginiamole come una serie di porte da oltrepassare, dove l'input � l'uomo che le deve oltrepassare e ogni volta che lo fa cambia qualcosa nel suo comportamento, finch� all'ultima porta oltrepassata, l'uomo � diventato una persona del tutto differente rappresentando l'output di questo processo."\\
Questa tesi si focalizza su un tipo particolare di algoritmo di machine learning: la \textit{Classification} (Classificazione), in particolar modo della classificazione di immagini. L'obiettivo principale � trovare un'architettura ottimale per l'algoritmo che identifica le immagini di oggetti architettonici, per fare questo bisogna trovare la giusta topologia, la giusta profondit� e la giusta larghezza di ogni livello della rete neurale.
\chapter{Algoritmi di apprendimento}
Gli algoritmi di machine learning sono solitamente divisi in tre tipi principali:
\begin{itemize}
\item Supervised learning (apprendimento supervisionato)
\item Unsupervised learning (apprendimento non supervisionato)
\item Reinforcement learning (apprendimento per rinforzo)
\end{itemize}


\section{Apprendimento supervisionato}
Vengono presentati i dati di input e i risultati desiderati. Lo scopo � di apprendere una regola generale che colleghi i dati in ingresso con quelli in uscita.

\subsection{Classificazione}
\section{Apprendimento non supervisionato}
\section{Apprendimento per rinforzo}
\chapter{Reti neurali}
\end{document}
